\section*{Hoofdstuk \ref{ch:g12:ph-conductivity-titrations} --- Titraties met pH en geleidbaarheid}

\begin{solution}{exo:g12:ph-conductivity-titrations:1}
$V_E = \qty{20.0}{mL}$; $c = 0.100 \times 20.0 / 20.0 = \qty{0.100}{mol/L}$.
\end{solution}

\begin{solution}{exo:g12:ph-conductivity-titrations:2}
Bij het \omterm{def:g12:ph-conductivity-titrations:ph-metric}{halfequivalentiepunt} is \omterm{def:g9:acids-bases-ph:ph}{pH} $= pK_a$: $pK_a = 3.75$, dicht bij
methaanzuur (3.74).
\end{solution}

\begin{solution}{exo:g12:ph-conductivity-titrations:3}
Fenolftaleïne: zijn \omterm{def:g12:buffers-predominance:indicator}{omslagtraject} 8.0--10.0 ligt in de sprong, die
basisch is. Methylrood (4.2--6.3) zou omslaan in de buurt van het
\omterm{def:g12:ph-conductivity-titrations:ph-metric}{halfequivalentiepunt}, veel te vroeg.
\end{solution}

\begin{solution}{exo:g12:ph-conductivity-titrations:4}
De respons van de elektrode verloopt met de tijd en de temperatuur: hij
wordt zo ingesteld dat hij goed afleest in twee \omterm{def:g12:buffers-predominance:buffer}{bufferoplossingen} met een
bekende \omterm{def:g9:acids-bases-ph:ph}{pH}.
\end{solution}

\begin{solution}{exo:g12:ph-conductivity-titrations:5}
\ce{H3O+} (349.8) > \ce{OH-} (198.3) > \ce{Cl-} (76.2) > \ce{Na+} (50.3).
\end{solution}

\begin{solution}{exo:g12:ph-conductivity-titrations:6}
$c = 0.050 \times 14.6 / 20.0 = \qty{0.0365}{mol/L}$.
\end{solution}

\begin{solution}{exo:g12:ph-conductivity-titrations:7}
Hellingen: 2.0, 4.5, 17.5 en \qty{2.5}{mL^{-1}}. De grootste ligt tussen
10.0 en \qty{10.2}{mL}: $V_E \approx \qty{10.1}{mL}$.
\end{solution}

\begin{solution}{exo:g12:ph-conductivity-titrations:8}
Vóór het \omterm{def:g11:titration:equivalence}{equivalentiepunt} haalt elk toegevoegd \ce{OH-} een \ce{H3O+}
weg (dat heel goed geleidt) en brengt het een \ce{Na+} binnen (dat veel
minder goed geleidt): de \omterm{def:g12:ph-conductivity-titrations:conductivity}{geleidbaarheid} daalt. Daarna hopen \ce{Na+} en
\ce{OH-} zich op zonder te reageren: ze stijgt.
\end{solution}

\begin{solution}{exo:g12:ph-conductivity-titrations:9}
Aan het begin weinig \omterm{def:g9:ions:ion}{ionen} (het \omterm{def:g12:ka-and-pka:strong-acid}{zwakke zuur} reageert nauwelijks met
water): lage \omterm{def:g12:ph-conductivity-titrations:conductivity}{geleidbaarheid}. Vóór het \omterm{def:g11:titration:equivalence}{equivalentiepunt} ontstaan
ethanoaat- en natriumionen in plaats van ongeladen \omterm{def:g7:atoms-and-molecules:molecule}{moleculen}: ze stijgt,
langzaam. Na het \omterm{def:g11:titration:equivalence}{equivalentiepunt} hopen \ce{Na+} en \ce{OH-} zich op: ze
stijgt sneller. Twee stijgende lijnen, de tweede steiler.
\end{solution}

\begin{solution}{exo:g12:ph-conductivity-titrations:10}
Ongeveer 2.9, 4.76 en 8.7. Ethaanzuur is zwak: maar een paar procent
ervan levert \omterm{def:g12:ka-and-pka:oxonium}{oxoniumionen}, dus de \omterm{def:g9:acids-bases-ph:ph}{pH} begint hoger dan de 1.0 van zoutzuur
van ongeveer dezelfde concentratie.
\end{solution}

\begin{solution}{exo:g12:ph-conductivity-titrations:11}
Zodat het toegevoegde volume \omterm{def:g11:titration:titration}{titrant} het totale volume nauwelijks
verandert: de concentraties, en dus de \omterm{def:g12:ph-conductivity-titrations:conductivity}{geleidbaarheid}, veranderen dan
langs rechte lijnen.
\end{solution}

\begin{solution}{exo:g12:ph-conductivity-titrations:12}
De \omterm{def:g9:ions:ion}{hydroxide-ionen} reageren eerst met het \omterm{def:g12:ka-and-pka:strong-acid}{sterke zuur} (\ce{H3O+}), daarna
met het zwakke. De eerste sprong markeert het eind van het zoutzuur (zijn
\omterm{def:g11:titration:equivalence}{equivalentievolume} meet het), de tweede het eind van het ethaanzuur (het
volume tussen de twee sprongen meet het).
\end{solution}

\begin{solution}{exo:g12:ph-conductivity-titrations:13}
Begin: $(349.8 + 76.2) \times 0.0100 = \qty{4.26}{mS/cm}$. Bij het
\omterm{def:g11:titration:equivalence}{equivalentiepunt}: $[\ce{Na+}] = [\ce{Cl-}] = 1.00/110 = \qty{0.00909}{mol/L}$,
dus $(50.3 + 76.2) \times 0.00909 = \qty{1.15}{mS/cm}$. Allebei zoals in
de figuur.
\end{solution}

\begin{solution}{exo:g12:ph-conductivity-titrations:14}
Geen in $V_E$: de sprong ligt bij hetzelfde volume, wat de verschuiving
van de aflezingen ook is. De $pK_a$ die bij het \omterm{def:g12:ph-conductivity-titrations:ph-metric}{halfequivalentiepunt} wordt
afgelezen, is 0.2 te hoog.
\end{solution}

\begin{solution}{exo:g12:ph-conductivity-titrations:15}
De \omterm{def:g9:acids-bases-ph:ph}{pH} verandert niet tijdens deze \omterm{def:g9:ions:precipitate}{neerslagreactie}, maar de \omterm{def:g9:ions:ion}{ionen} wel:
vóór het \omterm{def:g11:titration:equivalence}{equivalentiepunt} haalt elk toegevoegd \ce{Ag+} een \ce{Cl-} weg
en brengt het een
\ce{NO3-} met een vergelijkbare \omterm{def:g12:ph-conductivity-titrations:conductivity}{geleidbaarheid} binnen, dus de
\omterm{def:g12:ph-conductivity-titrations:conductivity}{geleidbaarheid} blijft vrijwel constant (de natriumionen waren er vanaf
het begin); na het \omterm{def:g11:titration:equivalence}{equivalentiepunt} hopen \ce{Ag+} en \ce{NO3-}
zich op en stijgt ze. Een vlakke lijn, en daarna een stijgende.
\end{solution}

\begin{solution}{pb:g12:ph-conductivity-titrations:1}
\textbf{1.} \qty{6}{g} ethaanzuur per \qty{100}{g} azijn.

\textbf{2.} \qty{101}{g} azijn bevat $6.0 \times 101/100 =
\qty{6.06}{g}$; $M = \qty{60.0}{g/mol}$: \qty{0.101}{mol} in
\qty{0.1000}{L}, dat is \qty{1.01}{mol/L}.

\textbf{3.} Onverdund zou \qty{10.0}{mL} ongeveer \qty{101}{mL} natronloog
van \qty{0.100}{mol/L} nodig hebben: meer dan er in een buret past.

\textbf{4.} Een volumepipet van \qty{10.0}{mL} en een maatkolf van
\qty{100.0}{mL}.

\textbf{5.} \ce{CH3COOH + OH- -> CH3COO- + H2O}.

\textbf{6.} $V_E = \qty{10.1}{mL}$.

\textbf{7.} $c = 0.100 \times 10.1 / 10.0 = \qty{0.101}{mol/L}$.

\textbf{8.} $10 \times 0.101 = \qty{1.01}{mol/L}$.

\textbf{9.} \omterm{def:g9:acids-bases-ph:ph}{pH} 4.76 bij \qty{5.05}{mL}: de $pK_a$ van ethaanzuur.

\textbf{10.} Bij het \omterm{def:g11:titration:equivalence}{equivalentiepunt} bevat de oplossing natriumethanoaat,
en ethanoaat is een \omterm{def:g12:ka-and-pka:strong-acid}{zwakke base}.

\textbf{11.} Fenolftaleïne, waarvan het \omterm{def:g12:buffers-predominance:indicator}{omslagtraject} 8.0--10.0 de \omterm{def:g9:acids-bases-ph:ph}{pH} bij
het \omterm{def:g11:titration:equivalence}{equivalentiepunt} (ongeveer 8.7) bevat.

\textbf{12.} Natriumionen en ethanoaationen.

\textbf{13.} Elk toegevoegd \ce{OH-} maakt van een vrijwel niet
geïoniseerd \omterm{def:g7:atoms-and-molecules:molecule}{molecuul} een
\ce{CH3COO-}, met een \ce{Na+} erbij: er komen langzaam \omterm{def:g9:ions:ion}{ionen} met een
bescheiden \omterm{def:g12:ph-conductivity-titrations:conductivity}{geleidbaarheid} bij.

\textbf{14.} Na het \omterm{def:g11:titration:equivalence}{equivalentiepunt} hopen de \omterm{def:g9:ions:ion}{hydroxide-ionen}, die heel
goed geleiden, zich op met de natriumionen.

\textbf{15.} Het \omterm{def:g12:ka-and-pka:strong-acid}{zwakke zuur} geeft maar heel weinig \omterm{def:g9:ions:ion}{ionen} af aan water.

\textbf{16.} Bij $V_E = \qty{10.1}{mL}$, net als bij de \omterm{def:g12:ph-conductivity-titrations:ph-metric}{pH-metrische
titratie}.

\textbf{17.} $1.01 \times 60.0 = \qty{60.6}{g/L}$.

\textbf{18.} \qty{6.06}{g} in \qty{100.0}{mL}.

\textbf{19.} $6.06 \times 100 / 101 = \qty{6.0}{g}$ per \qty{100}{g}.

\textbf{20.} De azijn heeft een zuurgraad van \qty{6.0}{\percent}: het etiket
klopt.
\end{solution}
